\section{Discussion}

The aggregate result changes meaning once it is decomposed. Colombia's production-account TFP was nearly flat over 2005--2024, but this was not the common outcome of all activities. Social services, trade, agriculture, and transport made positive contributions that were almost fully offset by finance and real estate, mining, construction, manufacturing, and electricity. The relevant empirical question is therefore not only why aggregate TFP stagnated. It is also why productivity gains persisted in some activities while negative contributions remained large in others.

The results do not identify the mechanisms behind each activity path. Mining illustrates the limitation. Its output grew by 1.70 percent per year, but the measured input contributions exceeded that growth and produced TFP of $-2.28$ points. This pattern may reflect resource depletion, the composition of extraction, investment cycles, capacity utilization, regulation, or measurement. The data used here cannot distinguish among these explanations.

Finance and real estate present a different problem. The combined aggregate covers financial intermediation, real estate, and business and rental activities. Its negative contribution may therefore combine distinct subactivities. Output and price measurement are also difficult in financial and real-estate services. A useful extension would reconstruct the contribution at a finer level before assigning a common mechanism to the aggregate.

Manufacturing and construction show why both the sectoral change and the aggregation weight matter. Manufacturing had a moderate TFP decline but a large average weight. Construction had a sharper decline and a smaller weight. Their average aggregate contributions were almost identical. A policy diagnosis based only on the largest sectoral decline would miss manufacturing's aggregate relevance.

The counterfactual provides a scale for these negative contributions, not a policy target. Raising sectoral TFP requires changes in technology, organization, regulation, factor allocation, capacity use, or measurement, depending on the activity. A credible policy exercise must replace the zero-TFP benchmark with attainable sector-specific references and allow labor, capital, intermediate inputs, prices, and activity weights to respond.

The next empirical step should therefore be sector-specific. For manufacturing, the evidence should test the roles of technology adoption, management practices, exporting, and firm reallocation. For construction, it should distinguish project-cycle and capacity-utilization effects from persistent weaknesses in project management and production methods. For mining, it should separate resource and composition effects from investment and regulatory conditions. For finance and real estate, the first requirement is a more detailed statistical decomposition.

The analysis also has two general limitations. First, the period begins in 2005 because the sectoral production-account history in the annex is reported as annual changes rather than a longer published index. Second, the recovered weights are implicit in the DANE aggregation. Their exact reconciliation supports the contribution calculation, but the public annex does not expose the underlying nominal series used to construct them. Replication therefore validates the accounting system, not every stage of DANE's source-data construction.
